\ifdefined\iscombined
	\lecturetitle{Viewpoints}
\else
\documentclass{article}
\usepackage{listings}
\usepackage{amsthm}
\usepackage{comment}
\usepackage{amsmath}
\usepackage{amsfonts}
\usepackage{sectsty}
\usepackage{hyperref}
\usepackage{fancyhdr}
\usepackage[top=1in,bottom=1in,left=1in,right=1in]{geometry}
\usepackage{float}
\usepackage{graphicx}
\usepackage{tabularx}
\usepackage{mathtools}

\date{
	\vspace{-.75em}
	{\large \today}
}

\title{
	\vspace*{-1.5em}
	{\Huge Lecture 6} \\
	\vspace{-1em}
	\rule{\textwidth/4}{.01em} \\
	\vspace{-.25em}
	{\Large Viewpoints}
	\vspace{-.75em}
}

\author{
	{\large Matthew Farah}
}

\hypersetup{
	colorlinks=true,
	linkcolor=blue,
	filecolor=magenta,
	urlcolor=blue,
	citecolor=blue,
	anchorcolor=blue
}

\fancyhead{}
\fancyhead[L]{\today}
\fancyhead[R]{Lecture 6 - Viewpoints}
\sectionfont{\fontsize{12}{12}\bfseries}
\subsectionfont{\fontsize{10}{12}\normalfont}
\subsubsectionfont{\fontsize{10}{12}\normalfont}

\begin{document}

\maketitle
\pagestyle{fancy}

\fi

\begin{itemize}
	\item When considering a written, non-technical description of a system:
	\begin{itemize}
		\item Verbs often describe the system's functions.
		\item Nouns often describe what the system's functions interacts with.
	\end{itemize}
	\item The different modes of a system represent an equivalence class of the system's state.
	\begin{itemize}
		\item A system's state describes its behaviour in some way.
		\item A system's mode changes in response to business events.
	\end{itemize}
	\item All viewpoints contribute to a system's global scenario.
	\item Designs should ideally be open for addition but closed for modification.
\end{itemize}

\ifdefined\iscombined
	\newpage
\else
	\end{document}
\fi
